\documentclass[a4paper,12pt]{article}

\usepackage[T1]{fontenc}
\usepackage[italian]{babel}
\usepackage{graphicx}
\usepackage{geometry}
\usepackage{tabularx}
\usepackage[table]{xcolor}
\usepackage[most]{tcolorbox}
\usepackage{fancyhdr}
\usepackage[hidelinks]{hyperref}

\newcommand{\groupmail}{alittlebyte.swe@gmail.com}
\newcommand{\grouplogo}{../../../.github/site-src/assets/logo_alittlebyte.png}
\newcommand{\groupsymbol}{../../../.github/site-src/assets/solo_logo.png}

\geometry{
    left=2.5cm,
    right=2.5cm,
    top=2.5cm,
    bottom=2.5cm,
    headheight=331pt,
    includeheadfoot
}

\pagestyle{fancy}
\fancyhf{}

\renewcommand{\headrulewidth}{0pt}
\fancyhead{}

\renewcommand{\footrulewidth}{0.4pt}
\fancyfoot[L]{\includegraphics[height=0.6cm]{\groupsymbol}\hspace{0.45em}aLittleByte}
    \fancyfoot[C]{\thepage}
\fancyfoot[R]{Verbale 2026-06-24}

\fancypagestyle{firstpage}{
    \fancyhf{}
    \renewcommand{\headrulewidth}{0pt}
    \renewcommand{\footrulewidth}{0.4pt}
    \fancyhead[C]{%
        \makebox[\textwidth][c]{%
            \begin{tabular}{c}
                \includegraphics[height=10cm]{\grouplogo}\\[1ex]
                \small\texttt{\groupmail}
            \end{tabular}%
        }
    }
    \fancyfoot[L]{\includegraphics[height=0.6cm]{\groupsymbol}\hspace{0.45em}aLittleByte}
        \fancyfoot[C]{\thepage}
    \fancyfoot[R]{Verbale 2026-06-24}
}

\providecommand{\glslink}[2]{\href{http://127.0.0.1:8765/glossary-html\#gls-#1}{\underline{#2}\textsuperscript{\scriptsize G}}}

\begin{document}

\thispagestyle{firstpage}

\vspace*{0.5cm}

\begin{center}
    {\LARGE \textbf{Verbale Interno - 2026-06-24}}
\end{center}

\vspace{1.5cm}

\begin{center}
    \renewcommand{\arraystretch}{1.3}
    \begin{tcolorbox}[
        width=0.92\textwidth,
        colback=gray!6,
        colframe=gray!45,
        boxrule=0.5pt,
        arc=2mm,
        boxsep=0pt,
        left=8pt, right=8pt, top=8pt, bottom=8pt
    ]
        \begin{tabularx}{\linewidth}{>{\bfseries}p{3.5cm} X}
            Gruppo      & aLittleByte (gruppo 19) \\
            Redattore   & Diego Belluz \\
            Verificatori & Alex Oloni\\
            Destinatari & Prof. Tullio Vardanega, Prof. Riccardo Cardin \\
            Data        & 2026-06-24\\
        \end{tabularx}
    \end{tcolorbox}
\end{center}

\newpage

\newgeometry{left=2.5cm,right=2.5cm,top=2.5cm,bottom=2.5cm,includefoot}
\tableofcontents
\newpage
\section{Informazioni Generali}
\section*{Specifiche Documento}
\hrule
\vspace{1cm}

\begin{tabular}{>{\bfseries}p{5cm} | p{10cm}}
Luogo Riunione & \glslink{discord}{Discord} \\[6pt]

Orario (Inizio - Fine) & 15:30 -- 16:00 \\[6pt]

\glslink{redattore}{Redattore} &
  D. Belluz\\[6pt]

\glslink{amministratore}{Amministratore} & 
A. Gastaldello\\[6pt]

\glslink{verificatore}{Verificatore} &
A. Oloni\\[6pt]

Partecipanti &
  A. Oloni\\ &
  L. Soligo\\&
  A. Gastaldello\\&
  E. Bellet\\&
  D. Belluz\\&
  A. Maule\\&
  V. Y. Kaneda Fernandes\\[6pt]

\end{tabular}


\section{Ordine del Giorno}
    \begin{itemize}
    \item Preparazione della candidatura \glslink{rtb-requirements-and-technology-baseline}{RTB} 
    \item Suddivisione dello studio delle tecnologie 
    \item Richiesta approvazione PoC dell'azienda
\end{itemize}


\section{Attività svolte}
\subsection{Preparazione della Candidatura RTB}
Il gruppo ha fatto una valutazione generale riguardo a tutti i documenti che sono stati sistemati, ed è stato accertato che bisogna semplicemente approvarli o controllare alcune piccole sezioni. È stato quindi deciso di presentare nuovamente
la nostra candidatura per la revisione \glslink{rtb-requirements-and-technology-baseline}{RTB} per il giorno Venerdi 26 Giugno. 

\subsection{Suddivisione della spiegazione delle tecnologie}
In vista del colloquio, il gruppo ha concordato la necessità di saper giustificare il motivo per cui
sono state adottate certe scelte tecniche. Le singole tecnologie impiegate nel prodotto software
sono state quindi ripartite equamente tra i membri del team. Questo permetterà uno studio
individuale mirato, in modo che ciascuno possa esporre e motivare una parte specifica durante
l’incontro.

\subsection{Richiesta dettagli all'azienda}
Abbiamo chiesto l'approvazione delle tecnologie, tramite Telegram, ossia il canale di messaggistica predefinito dall'azienda per questo progetto, che abbiamo sostituito nel PoC all'azienda \glslink{committente-proponente}{committente}, principalmente l'utilizzo di \glslink{angular}{Angular}, la quale ci ha dato esito positivo.

\section{To-Do List}
\begin{table}[h]
    \centering
    \begin{tabular}{|p{4.5cm}|p{5cm}|l|}
    \hline
        \rowcolor{lightgray}\textit{\textbf{Persona Incaricata}}  & \textbf{Task Assegnata} &\textbf{Link \glslink{issue}{Issue}} \\
    \hline
        \textit{[Tutti i membri del gruppo]} & Studio approfondito della propria tecnologia assegnata & \href{https://github.com/aLittleByte-19/Documentazione/issues/133}{\#133}\\
    \hline
        \textit{D. Belluz} & Stesura \glslink{verbale-interno}{Verbale Interno} 24/06/2026 & \href{https://github.com/aLittleByte-19/Documentazione/issues/134}{\#134}\\
    \hline
    \end{tabular}
    \label{tab:placeholder}
\end{table}

\end{document}
